\section{\gretl contributed packages}\label{sec:gretl_contributed_packages}
There\footnote{This section is based on \cite{gretl_function_package_guide}.}
are two types of function packages in \gretl: those
contributed and maintained by users, and those developed by
the \gretl team. The latter are more tightly integrated with
the \gretl program than standard user-contributed packages.
This second type is referred to as \textit{add-ons} and is
automatically bundled with each \gretl distribution
(as is also the case for \gretlpy). Users can browse
automatically updated listings of \gretl \textit{add-ons}
at \url{https://gretl.sourceforge.net/addons-data/addons.xml}
and of contributed function packages at
\url{https://gretl.sourceforge.net/cgi-bin/gretldata.cgi}.
In both cases, documentation for the packages is available.

The core component of a \gretl function package---in simpler
cases, the sole component---is a \texttt{gfn} file. This is an
XML file conforming to the Document Type Definition
\texttt{gretlfunc.dtd}, which is supplied with the \gretl
distribution; the latest version is always available online.
Such files contain
\begin{itemize}
	\item \texttt{hansl} code for at least one function;
	\item metadata such as author, version, and release date;
	\item help text for the function(s), or a reference to
		documentation in PDF format;
	\item a sample script illustrating calls to the packaged function(s).
\end{itemize}

An alternative to a simple package is what we refer to as a
``zip package''. Such packages take the form of a \texttt{PKZIP}
archive containing a \texttt{gfn} file along with additional
materials. These may include PDF documentation (in place of
plain-text help), data required by the package for internal
use (for example, tables of critical values for a test statistic),
and/or supplementary data files or scripts intended to complement
the mandatory sample script. The list of \textit{add-ons} is as follows:
\begin{itemize}
	\item \textbf{dbnomics}: Obtain data via DBnomics, written by
		Riccardo ``Jack'' Lucchetti and Allin Cottrell, JEL codes: C81.\newline
		Documentation: \url{https://sourceforge.net/projects/gretl/files/addons/doc/dbnomics.pdf}.	
	\item \textbf{extra}: Extra functions for \texttt{hansl} scripting,
		written by the \gretl team, JEL codes: C88.\newline
		Documentation: \url{https://sourceforge.net/projects/gretl/files/addons/doc/extra.pdf}.	
	\item \textbf{geoplot}: Geospatial plotting utilities,
		written by Allin Cottrell and Riccardo ``Jack'' Lucchetti,
		JEL codes: C81, C88.\newline
		Documentation: \url{https://sourceforge.net/projects/gretl/files/addons/doc/geoplot.pdf}.	
	\item \textbf{gig}: An assortment of univariate GARCH models,
		written by Riccardo ``Jack'' Lucchetti and Stefano Balietti,
		JEL codes: C22.\newline
		Documentation: \url{https://sourceforge.net/projects/gretl/files/addons/doc/gig.pdf}.	
	\item \textbf{HIP}: Heteroskedastic IV Probit models,
		written by Riccardo ``Jack'' Lucchetti and Claudia Pigini,
		JEL codes: C25, C26\newline
		Documentation: \url{https://sourceforge.net/projects/gretl/files/addons/doc/HIP.pdf}.	
	\item \textbf{ivpanel}: Within- and between-panel TSLS models,
		plus G2SLS, written by Allin Cottrell, JEL codes: C23, C26.\newline
		Documentation: \url{https://sourceforge.net/projects/gretl/files/addons/doc/ivpanel.pdf}.	
	\item \textbf{KFgui}: Graphical user interface for Kalman filter
		models, written by Riccardo ``Jack'' Lucchetti, JEL codes: C32.\newline
		Documentation: \url{https://sourceforge.net/projects/gretl/files/addons/doc/KFgui.txt}.	
	\item \textbf{logging}: Structured logging for script execution,
		written by Artur Tarassow, JEL codes: C81.\newline
		Documentation: \url{https://sourceforge.net/projects/gretl/files/addons/doc/logging.pdf}.	
	\item \textbf{regls}: Regularized least-squares estimation
		methods, written by Allin Cottrell, JEL codes: C20, C52.\newline
		Documentation: \url{https://sourceforge.net/projects/gretl/files/addons/doc/regls.pdf}.	
	\item \textbf{SVAR}: Structural vector autoregression models,
		written by Riccardo ``Jack'' Lucchetti and Sven Schreiber,
		JEL codes: C32.\newline
		Documentation: \url{https://sourceforge.net/projects/gretl/files/addons/doc/SVAR.pdf}.
\end{itemize}


\subsection{Package installation}
To install a package, the user should call the \pythcode{pkg()} function
(see \ref{func:gretl.pkg}). Basically, there are two sources (and thus two
ways) of installing user-contributed packages:
\begin{itemize}
	\item the SourceForge server: a full list of packages is available
	at \url{https://gretl.sourceforge.net/cgi-bin/gretldata.cgi}
	;
	\item a local file.
\end{itemize}
To install a package, the user should call either
\pythcode{pkg('install', 'package_name')}
or, in the case of a local installation,
\pythcode{pkg('install', '/path/to/package_name', local=True)}.

We've talked about installing packages, but what exactly does it mean
for a package to be installed? Basically, it means that the package
file(s) are placed in a standard location—the \texttt{functions}
subdirectory—where \lgretl should always be able to find them
automatically. In each case, we're talking about a subdirectory
named \texttt{functions}, the actual path to which
differs by platform (and, on MS Windows, by national language):
\begin{itemize}
	\item On an English-language Windows installation, a typical value is:\newline
		\verb|C:\Users\login\AppData\Roaming\gretl\functions|\newline
		where ``login'' is a placeholder for your username.
	\item On Linux, it is likely to be:\newline
		\verb|$HOME/.gretl/functions|.
	\item On macOS:\newline
		\verb|$HOME/Library/Application Support/gretl/functions|.
\end{itemize}
Here \verb|$HOME| denotes your ``home'' directory, which is
always defined on Linux and macOS. However, you don't need to guess
at these locations: in the \python console you can call
\pythcode{gretl.get_accessor('\$gretldir')} to get the respective
base directory on your system: append ``\texttt{functions}''
and you'll have the canonical package paths.

Having installed a package, \gretlpy (and thus \lgretl) will automatically
find it in these places. Well, there's one possible catch to look out for.
The \texttt{gfn} file for a package may be placed in the \texttt{functions}
directory itself, or in a subdirectory with the same name as the package;
that is, on one or other of the following patterns:
\begin{verbatim}
functions/mypkg.gfn
functions/mypkg/mypkg.gfn
\end{verbatim}
Why the difference? If a package includes PDF documentation, or any
other files besides the \texttt{gfn}, \textit{it must be in its own
	subdirectory} (the second pattern). But if the package consists of
a \texttt{gfn} file only, then it should go directly into the \texttt{functions}
directory (the first pattern).  Gretl should get this right when
installing a package for you but if you install a package by hand,
copying it to the \texttt{functions} directory yourself, you should
pay attention to this point.
